% GENERATED FILE. Edit canonical lessons, then run npm run generate:books.
% canonical-source-hash: fnv1a64:f8942957f3d3d3f5
% canonical-lessons: BN-C130-manibyag, BN-C130-poysa, BN-C130-kham, BN-C130-daktikit, BN-C130-putul

\chapter{Wallet, Coin, Envelope, Stamp, Doll}
\label{ch:bn-wallet-coin-envelope-stamp-doll}
\hlchaptermodality{130}

\begin{chapteropening}
\textbf{By the end of this chapter:} \emph{I can say a wallet, a coin, an envelope, a stamp and a doll.}

Complete the last lesson of chapter 130: I can say a wallet, a coin, an envelope, a stamp and a doll.
\end{chapteropening}

\section[\texorpdfstring{\bn{মানিব্যাগ} (mānibyāg)}{মানিব্যাগ (mānibyāg)}]{\bn{মানিব্যাগ} (mānibyāg) — a wallet}

\label{lesson:BN-C130-manibyag}



\begin{quote}
Before the new one: say the Bengali for an eraser, then the Bengali for a bag.
\end{quote}

\subsection*{You'll want to know: \bn{মানিব্যাগ}}
\textbf{\bn{মানিব্যাগ}} — \emph{mānibyāg} — \textquotedblleft{}a wallet\textquotedblright{}.

\begin{grammarlens}[title={What you've built}]
One more word for this chapter.
\end{grammarlens}

\subsection*{Your turn}
\begin{itemize}
\raggedright
  \item \emph{Say it:} \emph{mānibyāg}
  \item \emph{Say it:} \emph{mānibyāg}, once more
  \item \emph{Say it:} say \emph{mānibyāg}
  \item \emph{Recall:} say the Bengali for an eraser, then the Bengali for a bag, then say \emph{mānibyāg} again
\end{itemize}

\begin{tcolorbox}[breakable,colback=teal!4,colframe=teal!35!black,title={Before you move on}]
What does \bn{মানিব্যাগ} mean? (\textquotedblleft{}A wallet\textquotedblright{}.) Say it once more.
\end{tcolorbox}
\section[\texorpdfstring{\bn{পয়সা} (pôysā)}{পয়সা (pôysā)}]{\bn{পয়সা} (pôysā) — a coin}

\label{lesson:BN-C130-poysa}



\begin{quote}
Before the new one: say the Bengali for a bag, then the Bengali for a wallet.
\end{quote}

\subsection*{You'll want to know: \bn{পয়সা}}
\textbf{\bn{পয়সা}} — \emph{pôysā} — \textquotedblleft{}a coin\textquotedblright{}.

\begin{grammarlens}[title={What you've built}]
One more word for this chapter.
\end{grammarlens}

\subsection*{Your turn}
\begin{itemize}
\raggedright
  \item \emph{Say it:} \emph{pôysā}
  \item \emph{Say it:} \emph{pôysā}, once more
  \item \emph{Say it:} say \emph{pôysā}
  \item \emph{Recall:} say the Bengali for a bag, then the Bengali for a wallet, then say \emph{pôysā} again
\end{itemize}

\begin{tcolorbox}[breakable,colback=teal!4,colframe=teal!35!black,title={Before you move on}]
What does \bn{পয়সা} mean? (\textquotedblleft{}A coin\textquotedblright{}.) Say it once more.
\end{tcolorbox}
\section[\texorpdfstring{\bn{খাম} (khām)}{খাম (khām)}]{\bn{খাম} (khām) — an envelope}

\label{lesson:BN-C130-kham}



\begin{quote}
Before the new one: say the Bengali for a wallet, then the Bengali for a coin.
\end{quote}

\subsection*{You'll want to know: \bn{খাম}}
\textbf{\bn{খাম}} — \emph{khām} — \textquotedblleft{}an envelope\textquotedblright{}.

\begin{grammarlens}[title={What you've built}]
One more word for this chapter.
\end{grammarlens}

\subsection*{Your turn}
\begin{itemize}
\raggedright
  \item \emph{Say it:} \emph{khām}
  \item \emph{Say it:} \emph{khām}, once more
  \item \emph{Say it:} say \emph{khām}
  \item \emph{Recall:} say the Bengali for a wallet, then the Bengali for a coin, then say \emph{khām} again
\end{itemize}

\begin{tcolorbox}[breakable,colback=teal!4,colframe=teal!35!black,title={Before you move on}]
What does \bn{খাম} mean? (\textquotedblleft{}An envelope\textquotedblright{}.) Say it once more.
\end{tcolorbox}
\section[\texorpdfstring{\bn{ডাকটিকিট} (ḍākṭikiṭ)}{ডাকটিকিট (ḍākṭikiṭ)}]{\bn{ডাকটিকিট} (ḍākṭikiṭ) — a stamp}

\label{lesson:BN-C130-daktikit}



\begin{quote}
Before the new one: say the Bengali for a coin, then the Bengali for an envelope.
\end{quote}

\subsection*{You'll want to know: \bn{ডাকটিকিট}}
\textbf{\bn{ডাকটিকিট}} — \emph{ḍākṭikiṭ} — \textquotedblleft{}a stamp\textquotedblright{}.

\begin{grammarlens}[title={What you've built}]
One more word for this chapter.
\end{grammarlens}

\subsection*{Your turn}
\begin{itemize}
\raggedright
  \item \emph{Say it:} \emph{ḍākṭikiṭ}
  \item \emph{Say it:} \emph{ḍākṭikiṭ}, once more
  \item \emph{Say it:} say \emph{ḍākṭikiṭ}
  \item \emph{Recall:} say the Bengali for a coin, then the Bengali for an envelope, then say \emph{ḍākṭikiṭ} again
\end{itemize}

\begin{tcolorbox}[breakable,colback=teal!4,colframe=teal!35!black,title={Before you move on}]
What does \bn{ডাকটিকিট} mean? (\textquotedblleft{}A stamp\textquotedblright{}.) Say it once more.
\end{tcolorbox}
\section[\texorpdfstring{\bn{পুতুল} (putul)}{পুতুল (putul)}]{\bn{পুতুল} (putul) — a doll}

\label{lesson:BN-C130-putul}



\begin{quote}
Before the new one: say the Bengali for an envelope, then the Bengali for a stamp.
\end{quote}

\subsection*{You'll want to know: \bn{পুতুল}}
\textbf{\bn{পুতুল}} — \emph{putul} — \textquotedblleft{}a doll\textquotedblright{}.

\begin{grammarlens}[title={What you've built}]
One more word for this chapter.
\end{grammarlens}

\subsection*{Your turn}
\begin{itemize}
\raggedright
  \item \emph{Say it:} \emph{putul}
  \item \emph{Say it:} \emph{putul}, once more
  \item \emph{Say it:} say \emph{putul}
  \item \emph{Recall:} say the Bengali for an envelope, then the Bengali for a stamp, then say \emph{putul} again
\end{itemize}

\begin{tcolorbox}[breakable,colback=teal!4,colframe=teal!35!black,title={Before you move on}]
What does \bn{পুতুল} mean? (\textquotedblleft{}A doll\textquotedblright{}.) Say it once more.
\end{tcolorbox}
